\documentclass[11pt,letterpaper]{article}
\usepackage{fontspec}
\setmainfont{texgyrepagella}[Extension=.otf,UprightFont=*-regular,BoldFont=*-bold,ItalicFont=*-italic,BoldItalicFont=*-bolditalic]
\usepackage[bidi=default]{babel}
\babelprovide[import,main]{spanish}
\babelprovide[import]{hebrew}
\babelprovide[import]{ancientgreek}
\babelfont[hebrew]{rm}[Path=fonts/,BoldFont=NotoHebrewSerif-700.ttf]{NotoHebrewSerif-400.ttf}
\babelfont[ancientgreek]{rm}[Path=fonts/]{GFSPorson.otf}
\newcommand{\heb}[1]{\mbox{\foreignlanguage{hebrew}{#1}}}
\newcommand{\hebl}[1]{\foreignlanguage{hebrew}{#1}}
\newcommand{\SENOR}{SEÑOR}
\newcommand{\hebverso}[1]{\par\smallskip{\centering\begin{otherlanguage}{hebrew}\large #1\end{otherlanguage}\par}\smallskip}
\newcommand{\grc}[1]{\foreignlanguage{ancientgreek}{#1}}
\usepackage[margin=2.5cm,headheight=14pt]{geometry}
\usepackage{microtype,ragged2e,booktabs,graphicx,xcolor,csquotes,array}
\newcolumntype{P}[1]{>{\RaggedRight\arraybackslash}p{#1}}
\usepackage[hidelinks]{hyperref}\usepackage{xurl}\urlstyle{same}
\usepackage[most]{tcolorbox}
\usepackage{tikz}\usetikzlibrary{arrows.meta,positioning,calc,decorations.pathreplacing}
\usepackage{titlesec,fancyhdr,caption}
\definecolor{tierraroja}{RGB}{122,40,18}\definecolor{caferojizo}{RGB}{58,18,8}
\definecolor{mostaza}{RGB}{214,160,58}\definecolor{papiro}{RGB}{243,228,192}\definecolor{sepia}{RGB}{98,66,44}
\titleformat{\section}{\Large\bfseries\color{tierraroja}}{\thesection.}{0.6em}{}
\titleformat{\subsection}{\large\bfseries\color{caferojizo}}{\thesubsection}{0.6em}{}
\pagestyle{fancy}\fancyhf{}
\fancyhead[L]{\small\itshape\color{sepia}\leftmark}\fancyhead[R]{\small\color{sepia}\thepage}
\renewcommand{\sectionmark}[1]{\markboth{#1}{}}
\renewcommand{\headrulewidth}{0.4pt}
\setlength{\parindent}{1.27cm}\linespread{1.15}
\captionsetup{font=small,labelfont={bf,color=tierraroja}}
\setlocalecaption{spanish}{table}{Tabla}
\setlength{\emergencystretch}{2.5em}
% --- Versos bíblicos: cuadro de un solo color (texto RVR1960 textual) ---
\newtcolorbox{verso}[1]{enhanced,unbreakable,colback=papiro,colframe=tierraroja,boxrule=0.9pt,arc=2pt,
  left=10pt,right=10pt,top=6pt,bottom=6pt,before skip=10pt,after skip=6pt,fontupper=\itshape,
  after upper={\par\hfill{\upshape\bfseries\small\textcolor{tierraroja}{(#1, RVR1960)}}}}
% --- Apartado filológico (sin cuadro de color) ---
\newenvironment{filologia}[1]{\par\smallskip\noindent{\small\bfseries\color{tierraroja}Apartado filológico}\enspace{\small\color{sepia}#1}\par\nopagebreak\smallskip\small\noindent
  \renewcommand{\arraystretch}{1.25}\begin{tabular}{@{}P{3.6cm} P{2.1cm} P{9.9cm}@{}}\toprule
  \textbf{Término} & \textbf{Raíz} & \textbf{Sentido y comentario}\\\midrule}
 {\bottomrule\end{tabular}\par\medskip}
\newcommand{\precision}[1]{\par\smallskip{\small\color{sepia}\textbf{Precisión.} #1}\par\smallskip}
\newcommand{\refapa}[1]{\par\hangindent=1.27cm\hangafter=1\noindent #1\par\vspace{4pt}}
\newcommand{\destacar}[1]{\textbf{\textcolor{tierraroja}{#1}}}
\title{\color{tierraroja}\bfseries La gloria de Dios\\[4pt]\large\color{sepia}\mdseries Conocer a Dios en sus atributos\\[2pt]\normalsize Introducción}
\author{}\date{}
\begin{document}
\maketitle\thispagestyle{empty}

\begin{verso}{Jeremías 9:23–24}
Así dijo \SENOR{}: No se alabe el sabio en su sabiduría, ni en su valentía se alabe el valiente, ni el rico se alabe en sus riquezas. Mas alábese en esto el que se hubiere de alabar: en entenderme y conocerme, que yo soy \SENOR{}, que hago misericordia, juicio y justicia en la tierra; porque estas cosas quiero, dice \SENOR{}.
\end{verso}

\begin{filologia}{Jer 9:22–23 en el texto hebreo (TM)}
\heb{אַל־יִתְהַלֵּל} \emph{al-yithalel} & \heb{הלל} \emph{halal} & “No se gloríe”. Forma reflexiva (\emph{hitpael}): alabarse \emph{a sí mismo}, jactarse. Es la misma raíz de \emph{Hallelu-Yah} (“alabad a Yah”): la alabanza que el ser humano se dirige a sí mismo, Dios la reclama para sí.\\
\heb{חָכָם} \emph{jakam} & \heb{חכם} & “Sabio”; de la misma raíz que \heb{חָכְמָה} (\emph{jojmá}), sabiduría práctica para vivir.\\
\heb{הַגִּבּוֹר} \emph{ha-guibor} & \heb{גבר} & “El valiente, el guerrero”; fuerza que se impone.\\
\heb{עָשִׁיר} \emph{ashir} & \heb{עשר} & “Rico”; la seguridad que dan los bienes.\\
\heb{הַשְׂכֵּל} \emph{haskel} & \heb{שׂכל} & “Entender”: comprender con prudencia y actuar con acierto.\\
\heb{וְיָדֹעַ אוֹתִי} \emph{ve-yadóa otí} & \heb{ידע} & “Y conocerme”: conocimiento por relación y experiencia; el mismo verbo describe la intimidad conyugal (Gn 4:1).\\
\heb{חֶסֶד} \emph{jésed} & \heb{חסד} & “Misericordia”: amor leal y comprometido, propio del pacto.\\
\heb{מִשְׁפָּט} \emph{mishpat} & \heb{שׁפט} & “Juicio”: el derecho que Dios establece y ejecuta.\\
\heb{צְדָקָה} \emph{tsedaqá} & \heb{צדק} & “Justicia”: rectitud que se hace realidad en las relaciones.\\
\heb{חָפַצְתִּי} \emph{jafatsti} & \heb{חפץ} & “Quiero, me deleito”: no solo lo que Dios hace, sino lo que le agrada.\\
\end{filologia}

\section{Una gloria fuera de lugar}

Toda esta obra nace de un versículo. Jeremías habla a un pueblo que está a punto de perderlo todo: Judá confía en la sabiduría de sus consejeros, en la fuerza de sus alianzas militares y en la prosperidad de sus ciudades, mientras el juicio de Babilonia se acerca. En ese contexto, Dios desmonta las tres grandes seguridades humanas —el saber, el poder y el tener— y señala la única gloria que permanece: \destacar{entenderle y conocerle}.

El verbo que usa el profeta es revelador. \heb{יִתְהַלֵּל} (\emph{yithalel}) es la forma reflexiva de \heb{הלל} (\emph{halal}), la raíz de la alabanza. El ser humano fue creado para alabar, pero cuando no alaba a Dios, se alaba a sí mismo. Jeremías no pide que dejemos de gloriarnos; pide que cambiemos el \emph{objeto} de nuestra gloria. Y ese objeto no es una idea, sino una Persona que se describe a sí misma por lo que hace y lo que ama: \emph{jésed}, \emph{mishpat} y \emph{tsedaqá}. Ya en el versículo que inaugura este libro aparecen los atributos de Dios.

Arthur W. Pink abrió su clásica obra sobre los atributos con este mismo texto, y lo resumió así: “Un conocimiento espiritual y salvador de Dios es la mayor necesidad de toda criatura humana” (Pink, 2020, p. 4). Y añadió una frase que funciona como justificación de todo lo que sigue: “No se puede confiar, ni servir, ni adorar a un Dios desconocido” (Pink, 2020, p. 4).

\section{Entender y conocer}

Jeremías une dos verbos que no son sinónimos. \emph{Entender} (\heb{שׂכל}) apunta a la mente: comprender cómo actúa Dios en la historia, qué valora y cómo gobierna. \emph{Conocer} (\heb{ידע}) apunta a la relación: vivir en pacto con Él, experimentar su trato. Un conocimiento sin el otro queda incompleto. La teología sin relación se vuelve fría; la experiencia sin entendimiento se vuelve frágil.

\begin{verso}{Oseas 6:3}
Y conoceremos, y proseguiremos en conocer a \SENOR{}; como el alba está dispuesta su salida, y vendrá a nosotros como la lluvia, como la lluvia tardía y temprana a la tierra.
\end{verso}

\begin{filologia}{Os 6:3}
\heb{וְנֵדְעָה} \emph{ve-nedeá} & \heb{ידע} & “Y conoceremos”: forma que expresa resolución o deseo (“¡conozcamos!”).\\
\heb{נִרְדְּפָה} \emph{nirdefá} & \heb{רדף} & “Proseguiremos”: literalmente “persigamos, corramos tras”. Conocer a Dios es una búsqueda activa y perseverante.\\
\heb{שַׁחַר} \emph{shájar} & \heb{שׁחר} & “Alba”: su venida es tan segura como el amanecer.\\
\end{filologia}

El Nuevo Testamento lleva esta idea a su punto más alto: conocer a Dios no es solo un deber; es la vida misma.

\begin{verso}{Juan 17:3}
Y esta es la vida eterna: que te conozcan a ti, el único Dios verdadero, y a Jesucristo, a quien has enviado.
\end{verso}

\begin{filologia}{Jn 17:3}
\grc{γινώσκωσιν} \emph{ginóskosin} & \grc{γινώσκω} & “Que conozcan”. Presente de subjuntivo: un conocer continuo, que no termina. La Septuaginta usa \grc{γινώσκω} para traducir \heb{ידע}.\\
\grc{ἀληθινόν} \emph{alethinón} & \grc{ἀλήθεια} & “Verdadero”: lo real y auténtico, en contraste con lo aparente.\\
\grc{ἡ αἰώνιος ζωή} \emph{he aiónios zoé} & \grc{ζωή} & “La vida eterna”: no solo duración sin fin, sino la vida propia de la era de Dios.\\
\end{filologia}

\section{Dios se da a conocer}

El Dios de la Escritura no se deja descubrir por investigación; se da a conocer. Pink lo afirmó sin rodeos: “Tal Dios no puede ser descubierto mediante investigación. Sólo puede ser conocido cuando es revelado al corazón por el Espíritu Santo, a través de la Palabra” (Pink, 2020, p. 8). Esta revelación tiene, sin embargo, varios medios: la creación, que muestra su poder; la Escritura, que muestra su carácter; y, de manera plena y definitiva, Jesucristo.

\begin{verso}{Salmo 19:1}
Los cielos cuentan la gloria de Dios, y el firmamento anuncia la obra de sus manos.
\end{verso}

\begin{filologia}{Sal 19:2 (TM)}
\heb{מְסַפְּרִים} \emph{mesaperim} & \heb{ספר} & “Cuentan”: relatan, narran. La creación es un relato continuo.\\
\heb{כְּבוֹד־אֵל} \emph{kevod-El} & \heb{כבד} & “La gloria de Dios”: el peso, la dignidad visible de lo que Dios es.\\
\heb{הָרָקִיעַ} \emph{ha-raquía} & \heb{רקע} & “El firmamento”: lo extendido, como metal batido (Gn 1:6).\\
\end{filologia}

\begin{verso}{Romanos 1:20}
Porque las cosas invisibles de él, su eterno poder y deidad, se hacen claramente visibles desde la creación del mundo, siendo entendidas por medio de las cosas hechas, de modo que no tienen excusa.
\end{verso}

\begin{filologia}{Ro 1:20}
\grc{τὰ ἀόρατα} \emph{ta aórata} & \grc{ὁράω} & “Las cosas invisibles”, con alfa privativa.\\
\grc{καθορᾶται} \emph{kathorátai} & \grc{ὁράω} & “Se ven claramente”: juego de palabras con lo anterior; lo invisible se hace visible.\\
\grc{θειότης} \emph{theiótes} & \grc{θεῖος} & “Deidad”: la cualidad divina que la creación deja ver. Se distingue de \grc{θεότης} (Col 2:9), la deidad misma, que habita en Cristo.\\
\end{filologia}

\precision{la revelación general muestra \emph{que} Dios existe y es poderoso, y deja al ser humano sin excusa; pero sus atributos morales y su plan de redención se conocen por la revelación especial: la Escritura y Cristo.}

\section{El Nombre: la colección de sus atributos}

El pasaje que mejor explica qué es un atributo está en el Sinaí. Después del becerro de oro, Moisés hace la petición más audaz de la Escritura:

\begin{verso}{Éxodo 33:18–19}
Él entonces dijo: Te ruego que me muestres tu gloria. Y le respondió: Yo haré pasar todo mi bien delante de tu rostro, y proclamaré el nombre de \SENOR{} delante de ti; y tendré misericordia del que tendré misericordia, y seré clemente para con el que seré clemente.
\end{verso}

\begin{filologia}{Éx 33:18–19}
\heb{כְּבֹדֶךָ} \emph{kevodeja} & \heb{כבד} & “Tu gloria”. La raíz significa “ser pesado”: la gloria es el peso real de lo que Dios es.\\
\heb{טוּבִי} \emph{tuví} & \heb{טוב} & “Mi bien, mi bondad”: Dios responde mostrando su carácter.\\
\heb{וְקָרָאתִי בְשֵׁם} \emph{ve-qaratí ve-shem} & \heb{קרא} + \heb{שׁם} & “Proclamaré el nombre”. En hebreo el nombre (\emph{shem}) no es una etiqueta, sino la expresión de la identidad de la persona.\\
\heb{וְחַנֹּתִי} \emph{ve-janotí} & \heb{חנן} & “Seré clemente”: gracia libre, no merecida.\\
\heb{וְרִחַמְתִּי} \emph{ve-rijamtí} & \heb{רחם} & “Tendré misericordia”: compasión entrañable.\\
\end{filologia}

Moisés pide ver la gloria; Dios responde proclamando su Nombre. Pink recoge esta conexión con una frase que se ha vuelto el principio de esta obra: “El nombre es la colección de sus atributos” (Pink, 2020, p. 100). La proclamación llega en el capítulo siguiente:

\begin{verso}{Éxodo 34:6–7}
Y pasando \SENOR{} por delante de él, proclamó: ¡\SENOR{}! ¡\SENOR{}! fuerte, misericordioso y piadoso; tardo para la ira, y grande en misericordia y verdad; que guarda misericordia a millares, que perdona la iniquidad, la rebelión y el pecado, y que de ningún modo tendrá por inocente al malvado; que visita la iniquidad de los padres sobre los hijos y sobre los hijos de los hijos, hasta la tercera y cuarta generación.
\end{verso}

\begin{filologia}{Éx 34:6–7}
\heb{אֵל} \emph{El} & \heb{אל} & “Fuerte”: Dios en su poder.\\
\heb{רַחוּם} \emph{rajum} & \heb{רחם} & “Misericordioso”. De \heb{רֶחֶם} (\emph{réjem}), “seno materno”: amor entrañable.\\
\heb{חַנּוּן} \emph{janún} & \heb{חנן} & “Piadoso, clemente”: el que concede favor inmerecido (gracia).\\
\heb{אֶרֶךְ אַפַּיִם} \emph{érej apáyim} & \heb{ארך} + \heb{אף} & “Tardo para la ira”; literalmente “largo de narices”: su enojo tarda en encenderse (paciencia).\\
\heb{רַב־חֶסֶד} \emph{rav-jésed} & \heb{חסד} & “Grande en misericordia”: abundante en amor leal.\\
\heb{וֶאֱמֶת} \emph{ve-emet} & \heb{אמן} & “Y verdad”: firmeza, confiabilidad; de la misma raíz que “amén”.\\
\heb{נֹשֵׂא עָוֹן} \emph{nosé avón} & \heb{נשׂא} & “Que perdona la iniquidad”; literalmente “que carga” con ella.\\
\heb{וְנַקֵּה לֹא יְנַקֶּה} \emph{ve-naqé lo yenaqé} & \heb{נקה} & “De ningún modo tendrá por inocente”: la misericordia no anula la justicia.\\
\end{filologia}

Este texto es el corazón del criterio de esta obra. Dios no se define con categorías abstractas, sino \destacar{proclamando su Nombre}, y en ese Nombre aparecen juntos su poder, su misericordia, su gracia, su paciencia, su amor leal, su verdad y su justicia. Es, además, el texto del Antiguo Testamento más citado dentro del propio Antiguo Testamento (Nm 14:18; Neh 9:17; Sal 86:15; 103:8; Jl 2:13; Jon 4:2).

\section{Los nombres que revelan los atributos}

Si el Nombre es la colección de los atributos, cada nombre con que Dios se da a conocer en la Escritura revela algo de lo que él es. Muchos de ellos señalan un atributo concreto: Dios se llama ``el Santo'', ``Dios Todopoderoso'', ``Dios de verdad'', ``Dios misericordioso''. No todos los atributos tienen un nombre propio, pero la gran mayoría está ligada a un nombre, a un título o a una declaración hecha con el Nombre. Por eso cada capítulo de este libro incluye una sección sobre \destacar{el Nombre que revela el atributo}.

La tradición judía expresó la misma idea en un midrás sobre la zarza ardiente. Moisés había preguntado: ``Si ellos me preguntaren: ¿Cuál es su nombre?, ¿qué les responderé?'' (Éx 3:13), y el midrás pone en boca de Dios esta respuesta:

\begin{quote}\small
\hebl{שמי אתה מבקש לידע לפי מעשי אני נקרא}\par\smallskip
``¿Mi nombre quieres saber? Según mis obras soy llamado. Cuando juzgo a las criaturas, soy llamado \emph{Elohim}; cuando hago guerra contra los impíos, soy llamado \emph{Tsevaot}; cuando suspendo el castigo por los pecados del hombre, soy llamado \emph{El Shadday}; y cuando tengo misericordia de mi mundo, soy llamado \SENOR{}'' (\emph{Shemot Rabá} 3:6).
\end{quote}

\precision{es una lectura rabínica, atribuida a Rabí Abbá bar Memel, y no una etimología de cada nombre. Se cita porque muestra que, para Israel, los nombres de Dios no eran sinónimos intercambiables: cada uno revelaba un aspecto de su obrar.}

\subsection{``YO SOY EL QUE SOY''}

El primer ejemplo es el Nombre revelado en la zarza:

\hebverso{וַיֹּאמֶר אֱלֹהִים אֶל־מֹשֶׁה אֶהְיֶה אֲשֶׁר אֶהְיֶה}
\begin{verso}{Éxodo 3:14–15}
Y respondió Dios a Moisés: YO SOY EL QUE SOY. Y dijo: Así dirás a los hijos de Israel: YO SOY me envió a vosotros. Además dijo Dios a Moisés: Así dirás a los hijos de Israel: \SENOR{}, el Dios de vuestros padres, el Dios de Abraham, Dios de Isaac y Dios de Jacob, me ha enviado a vosotros. Este es mi nombre para siempre; con él se me recordará por todos los siglos.
\end{verso}

\begin{filologia}{Éx 3:14–15}
\heb{אֶהְיֶה אֲשֶׁר אֶהְיֶה} \emph{Ehyeh asher ehyeh} & \heb{היה} & “Soy el que soy” o “seré el que seré”. \emph{Ehyeh} es el imperfecto de \emph{hayá}, “ser, estar, llegar a ser”: una acción no terminada. Dios es sin principio y sin agotarse.\\
\heb{יְהוָה} \emph{YHWH} & \heb{היה} (discutida) & El Nombre propio (v. 15). El texto lo relaciona con \emph{ehyeh}: Dios dice “Yo soy”; Israel dice “Él es”. La etimología exacta es discutida (Koehler y Baumgartner, 2001).\\
\heb{לְעֹלָם} \emph{le'olam} & \heb{עלם} & “Para siempre”. El Nombre no cambia con las generaciones: “con él se me recordará por todos los siglos” (\heb{לְדֹר דֹּר}, “de generación en generación”).\\
\grc{ἐγώ εἰμι ὁ ὤν} & --- & La Septuaginta traduce “Yo soy el que es”: subraya el ser eterno de Dios.\\
\end{filologia}

Este Nombre habla, ante todo, de la \destacar{eternidad} de Dios: él simplemente es, sin comienzo, sin dependencia y sin fin. Por eso será el texto del capítulo sobre el Dios eterno. Jesús lo tomó para sí: “Antes que Abraham fuese, yo soy” (\grc{ἐγὼ εἰμί}, Jn 8:58), y el Apocalipsis lo despliega en tres tiempos: “el que es y que era y que ha de venir” (\grc{ὁ ὢν καὶ ὁ ἦν καὶ ὁ ἐρχόμενος}, Ap 1:8).

\precision{Rashi, siguiendo el Talmud (Berajot 9b), lee \emph{Ehyeh asher ehyeh} como una promesa de presencia: “estaré con ellos en esta angustia”, como estaré con ellos en las que vendrán. La lectura judía subraya la presencia fiel; la Septuaginta y Juan 8:58, el ser eterno. No se excluyen: el que es siempre, está siempre con los suyos.}

\subsection{Mapa de los nombres}

La Tabla~\ref{tab:nombres} vincula cada atributo de este libro con el nombre, título o declaración que lo revela. Donde no hay un nombre propio, se indica la declaración hecha con el Nombre.

\begin{table}[htbp]
\centering\footnotesize
\caption{Los atributos y los nombres que los revelan}\label{tab:nombres}
\renewcommand{\arraystretch}{1.18}
\begin{tabular}{@{}P{2.3cm} P{5.0cm} P{4.6cm} P{3.4cm}@{}}
\toprule
\textbf{Atributo} & \textbf{Nombre, título o declaración} & \textbf{Hebreo / griego} & \textbf{Texto} \\ \midrule
\multicolumn{4}{@{}l}{\textcolor{tierraroja}{\textbf{I. Ser}}}\\
Preexistente & “Yo soy el primero, y yo soy el postrero”; “el Alfa y la Omega” & \heb{אֲנִי רִאשׁוֹן וַאֲנִי אַחֲרוֹן}; \grc{τὸ ἄλφα καὶ τὸ ὦ} & Is 44:6; Ap 1:8 \\
Uno & “El \SENOR{} uno es”; “uno su nombre” & \heb{יְהוָה אֶחָד} & Dt 6:4; Zac 14:9 \\
Amor & “El Dios de paz y de amor” & \grc{ὁ θεὸς τῆς ἀγάπης} & 2 Co 13:11; 1 Jn 4:8 \\
Espíritu & “Dios de los espíritus de toda carne” & \heb{אֵל אֱלֹהֵי הָרוּחֹת} & Nm 16:22; Jn 4:24 \\
\multicolumn{4}{@{}l}{\textcolor{tierraroja}{\textbf{II. Identidad}}}\\
Santo & “El Santo de Israel”; “cuyo nombre es el Santo” & \heb{קְדוֹשׁ יִשְׂרָאֵל}; \heb{קָדוֹשׁ שְׁמוֹ} & Is 41:14; 57:15 \\
Eterno & “YO SOY EL QUE SOY”; “Dios eterno” & \heb{אֶהְיֶה אֲשֶׁר אֶהְיֶה}; \heb{אֵל עוֹלָם} & Éx 3:14; Gn 21:33 \\
Soberano & “Dios Altísimo”; “el solo Soberano” & \heb{אֵל עֶלְיוֹן}; \grc{ὁ μόνος δυνάστης} & Gn 14:19; 1 Ti 6:15 \\
Justo & “\SENOR{}, justicia nuestra” & \heb{יְהוָה צִדְקֵנוּ} & Jer 23:6 \\
Inmutable & “Yo \SENOR{} no cambio”; “la Roca” & \heb{לֹא שָׁנִיתִי}; \heb{הַצּוּר} & Mal 3:6; Dt 32:4 \\
Omnipotente & “Dios Todopoderoso”; “\SENOR{} de los ejércitos” & \heb{אֵל שַׁדַּי}; \heb{יְהוָה צְבָאוֹת}; \grc{παντοκράτωρ} & Gn 17:1; Is 6:3; 2 Co 6:18 \\
Omnipresente & “\SENOR{}-sama” (“el \SENOR{} está allí”); “Emanuel” & \heb{יְהוָה שָׁמָּה}; \heb{עִמָּנוּאֵל} & Ez 48:35; Is 7:14 \\
Omnisciente & “Dios que ve”; “Dios de todo saber” & \heb{אֵל רֳאִי}; \heb{אֵל דֵּעוֹת} & Gn 16:13; 1 S 2:3 \\
\multicolumn{4}{@{}l}{\textcolor{tierraroja}{\textbf{III. Relación con el ser humano}}}\\
Sabiduría & “Admirable, Consejero”; “el único y sabio Dios” & \heb{פֶּלֶא יוֹעֵץ}; \grc{μόνῳ σοφῷ θεῷ} & Is 9:6; Ro 16:27 \\
Verdadero & “Dios de verdad”; “el Dios del Amén” & \heb{אֵל אֱמֶת}; \heb{אֱלֹהֵי אָמֵן} & Sal 31:5; Is 65:16; Jer 10:10 \\
Misericordioso & “Dios misericordioso”; “Padre de misericordias” & \heb{אֵל רַחוּם}; \grc{ὁ πατὴρ τῶν οἰκτιρμῶν} & Dt 4:31; 2 Co 1:3 \\
Fiel & “Dios fiel” & \heb{הָאֵל הַנֶּאֱמָן}; \grc{πιστὸς ὁ θεός} & Dt 7:9; 1 Co 1:9 \\
Celo & “Cuyo nombre es Celoso” & \heb{קַנָּא שְׁמוֹ} & Éx 34:14 \\
Ira santa & “Dios de las venganzas” & \heb{אֵל נְקָמוֹת} & Sal 94:1; Nah 1:2 \\
Bondad & “\SENOR{} proveerá”; “es bueno el \SENOR{}” & \heb{יְהוָה יִרְאֶה}; \heb{טוֹב יְהוָה} & Gn 22:14; Sal 34:8 \\
Paciencia & “Tardo para la ira”; “el Dios de la paciencia” & \heb{אֶרֶךְ אַפַּיִם}; \grc{ὁ θεὸς τῆς ὑπομονῆς} & Éx 34:6; Ro 15:5 \\
Gracia & “Piadoso”; “el Dios de toda gracia” & \heb{חַנּוּן}; \grc{ὁ θεὸς πάσης χάριτος} & Éx 34:6; 1 P 5:10 \\ \bottomrule
\end{tabular}
\par\smallskip{\footnotesize\color{sepia}\textit{Nota.} Elaboración propia. Las citas siguen la RVR1960; Sal 31:5 es 31:6 en hebreo e Is 9:6 es 9:5.}
\end{table}

\precision{la tabla distingue grados. Algunos son nombres propios o títulos (\emph{El Shadday}, \emph{El Elyon}, “el Santo de Israel”); otros son declaraciones con el Nombre (“Yo \SENOR{} no cambio”, Mal 3:6); otros, nombres de lugares que revelan a Dios (“\SENOR{} proveerá”, Gn 22:14, nombra la provisión antes que la bondad en general; “\SENOR{}-sama”, Ez 48:35, es el nombre de la ciudad). En Ro 15:5, \grc{ὑπομονή} es “constancia”; la longanimidad de Dios se expresa sobre todo con \grc{μακροθυμία} (Ro 2:4). La tradición rabínica añadió sus propios nombres: \emph{HaQadosh Barukh Hu}, “el Santo, bendito sea”; \emph{HaRajamán}, “el Misericordioso”; \emph{Ribbonó shel olam}, “Señor del mundo”; y \emph{HaMaqom}, “el Lugar”, porque Dios es el lugar del mundo y no el mundo su lugar (\emph{Bereshit Rabá} 68:9).}

\subsection{El Nombre que es sobre todo nombre}

Pablo, judío formado en el \emph{Shemá}, dio el último paso. Tomó un texto en que el \SENOR{} reclama para sí solo la adoración de toda la tierra y lo aplicó a Jesús:

\begin{verso}{Isaías 45:23}
Por mí mismo hice juramento, de mi boca salió palabra en justicia, y no será revocada: Que a mí se doblará toda rodilla, y jurará toda lengua.
\end{verso}

\begin{verso}{Filipenses 2:9–11}
Por lo cual Dios también le exaltó hasta lo sumo, y le dio un nombre que es sobre todo nombre, para que en el nombre de Jesús se doble toda rodilla de los que están en los cielos, y en la tierra, y debajo de la tierra; y toda lengua confiese que Jesucristo es el Señor, para gloria de Dios Padre.
\end{verso}

\begin{filologia}{Flp 2:9–11}
\grc{τὸ ὄνομα τὸ ὑπὲρ πᾶν ὄνομα} & --- & “El nombre que está por encima de todo nombre”.\\
\grc{κάμψῃ … ἐξομολογήσηται} & --- & “Se doble … confiese”: los mismos verbos de Is 45:23 en la Septuaginta (\grc{κάμψει … ἐξομολογήσεται}).\\
\grc{κύριος} \emph{kýrios} & --- & “Señor”: la palabra con que la Septuaginta traduce el Nombre divino \heb{יְהוָה}. Confesar que Jesús es \emph{kýrios} es reconocerle el Nombre.\\
\end{filologia}

Cada nombre de la Tabla~\ref{tab:nombres} converge en ese Nombre. Por eso cada capítulo termina en Cristo: en él habita corporalmente toda la plenitud de la Deidad (Col 2:9), y en él los nombres de Dios se ven con rostro humano (2 Co 4:6).

\section{¿Qué es un atributo?}

Un \destacar{atributo} es una perfección eterna que Dios revela de sí mismo en la Escritura y por la cual podemos conocerle. No es una parte de Dios, como si Él estuviera compuesto de piezas, sino Dios mismo considerado desde uno de los aspectos de su ser. La teología clásica llama a esto \emph{simplicidad divina}: Dios no \emph{tiene} amor y santidad como quien tiene objetos; Dios \emph{es} amor y \emph{es} santo, y lo es completamente a la vez (Berkhof, 1938).

De esta definición se desprenden tres principios que guiarán todo el libro. El primero es que \destacar{todos los atributos son eternos}: Dios no adquiere perfecciones con el tiempo ni las pierde.

\begin{verso}{Malaquías 3:6}
Porque yo \SENOR{} no cambio; por esto, hijos de Jacob, no habéis sido consumidos.
\end{verso}

\begin{filologia}{Mal 3:6}
\heb{לֹא שָׁנִיתִי} \emph{lo shaniti} & \heb{שׁנה} & “No cambio”. La raíz significa “cambiar, alterar” y también “repetir”. Dios no se altera ni en su ser ni en su carácter.\\
\heb{לֹא כְלִיתֶם} \emph{lo jelitem} & \heb{כלה} & “No habéis sido consumidos”: la inmutabilidad de Dios es la razón de la supervivencia de su pueblo.\\
\end{filologia}

El segundo principio es que \destacar{los atributos no se oponen}. La justicia de Dios no compite con su misericordia, ni su santidad con su amor. En Dios actúan en perfecta armonía:

\begin{verso}{Salmo 85:10}
La misericordia y la verdad se encontraron; la justicia y la paz se besaron.
\end{verso}

\begin{filologia}{Sal 85:11 (TM)}
\heb{נִפְגָּשׁוּ} \emph{nifgashu} & \heb{פגשׁ} & “Se encontraron”: salieron al encuentro una de la otra.\\
\heb{נָשָׁקוּ} \emph{nashaqu} & \heb{נשׁק} & “Se besaron”: imagen de reconciliación y unión.\\
\heb{חֶסֶד וֶאֱמֶת} \emph{jésed ve-emet} & — & El mismo binomio de Éx 34:6: amor leal y verdad.\\
\end{filologia}

El tercer principio es que \destacar{los atributos no se agotan}. Lo que conocemos de Dios es verdadero, pero no exhaustivo:

\begin{verso}{Job 26:14}
He aquí, estas cosas son sólo los bordes de sus caminos; ¡y cuán leve es el susurro que hemos oído de él! Pero el trueno de su poder, ¿quién lo puede comprender?
\end{verso}

\begin{filologia}{Job 26:14}
\heb{קְצוֹת} \emph{qetsot} & \heb{קצה} & “Bordes, extremos”: solo vemos la orilla de sus caminos.\\
\heb{שֵׁמֶץ} \emph{shémets} & \heb{שׁמץ} & “Susurro”: apenas un murmullo de lo que Dios es.\\
\heb{רַעַם} \emph{raam} & \heb{רעם} & “Trueno”: contraste entre lo que oímos (un susurro) y lo que Él es (un trueno).\\
\end{filologia}

\section{Primera clasificación: Ser, Identidad y Relación}

Toda clasificación de los atributos es una herramienta didáctica; la Escritura no presenta una lista sistemática. Este libro los organiza según tres preguntas, siguiendo un criterio hebreo: Dios se da a conocer por lo que \emph{es}, por \emph{quién es} y por \emph{cómo se relaciona} con el ser humano.

\subsection{Atributos propios de su Ser: ¿qué es Dios?}

Son las perfecciones que describen la naturaleza misma de Dios, eterna e independiente de toda creación. La Escritura las expresa con la fórmula más directa posible: “Dios es…”. Este libro reconoce tres: Dios es \emph{uno}, Dios es \emph{Espíritu} y Dios es \emph{amor}.

\begin{verso}{Deuteronomio 6:4}
Oye, Israel: \SENOR{} nuestro Dios, \SENOR{} uno es.
\end{verso}

\begin{filologia}{Dt 6:4}
\heb{שְׁמַע} \emph{shemá} & \heb{שׁמע} & “Oye”: escuchar para obedecer. Da nombre a la confesión central de Israel, el \emph{Shemá}.\\
\heb{יְהוָה} \emph{YHWH} & \heb{היה} & El Nombre personal de Dios, vinculado con el verbo “ser” (Éx 3:14). RVR1960 lo traduce “\SENOR{}”.\\
\heb{אֶחָד} \emph{ejad} & \heb{אחד} & “Uno”: unidad indivisible. El mismo término describe a varón y mujer como “una sola carne” (Gn 2:24).\\
\end{filologia}

\precision{que \heb{אֶחָד} pueda designar una unidad compuesta (Gn 2:24) es una observación frecuente en la teología cristiana; el término por sí mismo no demuestra la Trinidad, pero es compatible con ella.}

\begin{verso}{Juan 4:24}
Dios es Espíritu; y los que le adoran, en espíritu y en verdad es necesario que adoren.
\end{verso}

\begin{filologia}{Jn 4:24}
\grc{πνεῦμα ὁ θεός} \emph{pneuma ho theós} & \grc{πνέω} & “Dios es Espíritu”. \grc{πνεῦμα} viene de “soplar”; en el Antiguo Testamento su equivalente es \heb{רוּחַ} (\emph{rúaj}): viento, aliento, espíritu (Gn 1:2). Dios no es material ni está limitado por un cuerpo.\\
\end{filologia}

\begin{verso}{1 Juan 4:8}
El que no ama, no ha conocido a Dios; porque Dios es amor.
\end{verso}

\begin{filologia}{1 Jn 4:8}
\grc{ὁ θεὸς ἀγάπη ἐστίν} \emph{ho theós agápe estín} & \grc{ἀγαπάω} & “Dios es amor”. \grc{ἀγάπη} va sin artículo: no dice que el amor sea Dios, sino que el amor define la naturaleza de Dios.\\
\grc{οὐκ ἔγνω} \emph{ouk égno} & \grc{γινώσκω} & “No ha conocido”: el mismo verbo de Jn 17:3. Quien conoce a Dios, ama.\\
\end{filologia}

\precision{en el esquema de esta obra, el amor se nombra con \heb{חֶסֶד} (\emph{jésed}), el amor leal de pacto. Conviene notar que la Septuaginta suele traducir \emph{jésed} con \grc{ἔλεος} (“misericordia”), mientras que \grc{ἀγάπη} corresponde mejor a \heb{אַהֲבָה} (\emph{ahavá}; Dt 7:8). Ambos términos se estudiarán en el capítulo sobre el amor.}

\subsection{Atributos propios de su Identidad: ¿quién es Dios?}

Son las perfecciones que expresan el carácter propio y absoluto de Dios, aquello que lo distingue como el único Dios: su santidad, su eternidad, su soberanía, su justicia, su inmutabilidad, su omnipotencia, su omnipresencia y su omnisciencia. En el centro de todas ellas está la santidad, el único atributo que la Escritura repite tres veces:

\begin{verso}{Isaías 6:3}
Y el uno al otro daba voces, diciendo: Santo, santo, santo, \SENOR{} de los ejércitos; toda la tierra está llena de su gloria.
\end{verso}

\begin{filologia}{Is 6:3}
\heb{קָדוֹשׁ} \emph{qadosh} & \heb{קדשׁ} & “Santo”: apartado, distinto, absolutamente otro; y, a la vez, moralmente puro. La triple repetición es la forma hebrea del superlativo máximo.\\
\heb{צְבָאוֹת} \emph{tsevaot} & \heb{צבא} & “De los ejércitos”: Señor de las huestes celestiales.\\
\heb{מְלֹא כָל־הָאָרֶץ} \emph{meló kol-haárets} & \heb{מלא} & “Toda la tierra está llena”: base hebrea de la omnipresencia.\\
\heb{כְּבוֹדוֹ} \emph{kevodó} & \heb{כבד} & “Su gloria”: la santidad es lo que Dios es; la gloria es esa santidad hecha visible.\\
\end{filologia}

\subsection{Atributos de relación con el ser humano: ¿cómo se revela Dios?}

Aquí se requiere precisión. \destacar{Todos los atributos son perfecciones eternas de Dios}; ninguno comenzó a existir con la creación. Sin embargo, algunos de ellos \emph{se fueron revelando} en la historia de la relación de Dios con el ser humano. La sabiduría se hizo visible en la creación (Pr 3:19); la ira y la misericordia, ante el pecado y la miseria (Gn 3; Ro 1:18); la fidelidad, la verdad y el celo, en el pacto (Éx 20:5; 34:6); y todos ellos, en plenitud, en Cristo (Jn 1:14). Son la sabiduría, la verdad, la misericordia, la fidelidad, el celo y la ira santa; a ellos se suman la bondad, la paciencia y la gracia, que Pink estudia por separado y que el propio Éxodo 34:6 nombra.

\precision{decir que la ira o la misericordia son “de relación” no significa que sean temporales. Son la santidad y el amor eternos de Dios manifestados ante el pecado y la necesidad. Lo histórico es su \emph{manifestación}, no su existencia.}

La Tabla~\ref{tab:mapa} reúne los atributos que estudiará este libro, con su término hebreo y su texto ancla, y la Figura~\ref{fig:gloria} muestra cómo se relacionan entre sí.

\begin{table}[htbp]
\centering\small
\caption{Mapa de los atributos según la primera clasificación}\label{tab:mapa}
\renewcommand{\arraystretch}{1.12}
\begin{tabular}{@{}P{2.4cm} P{2.7cm} P{2.5cm} P{4.3cm} P{2.6cm}@{}}
\toprule
\textbf{Atributo} & \textbf{Hebreo} & \textbf{Transliteración} & \textbf{Rol} & \textbf{Texto ancla} \\ \midrule
\multicolumn{5}{@{}l}{\textcolor{tierraroja}{\textbf{I. Ser — ¿Qué es?}}}\\
Uno & \heb{אֶחָד} & \emph{ejad} & Base: unidad indivisible & Dt 6:4 \\
Amor & \heb{חֶסֶד} & \emph{jésed} & Eje: amor leal de pacto & Éx 34:6; 1 Jn 4:8 \\
Espíritu & \heb{רוּחַ} & \emph{rúaj} & Esencia no material & Gn 1:2; Jn 4:24 \\
\multicolumn{5}{@{}l}{\textcolor{tierraroja}{\textbf{II. Identidad — ¿Quién es?}}}\\
Santo & \heb{קָדוֹשׁ} & \emph{qadosh} & Centro que permea todo & Is 6:3 \\
Eterno & \heb{עוֹלָם} & \emph{olam} & Sin principio ni fin & Sal 90:2 \\
Soberano & \heb{אֲדֹנָי} & \emph{Adonay} & Gobierno supremo & Sal 115:3 \\
Justo & \heb{צֶדֶק} & \emph{tsedeq} & Norma de rectitud & Sal 97:2 \\
Inmutable & \heb{לֹא שָׁנִיתִי} & \emph{lo shaniti} & No varía & Mal 3:6 \\
Omnipotente & \heb{אֵל שַׁדַּי} & \emph{El Shadday} & Poder suficiente & Gn 17:1 \\
Omnipresente & \heb{מְלֹא כָל־הָאָרֶץ} & \emph{meló kol-haárets} & Presencia total & Is 6:3; Sal 139:7 \\
Omnisciente & \heb{יָדַע} & \emph{yadá} & Conocimiento exhaustivo & Sal 139:1–4 \\
\multicolumn{5}{@{}l}{\textcolor{tierraroja}{\textbf{III. Relación con el ser humano — ¿Cómo se revela?}}}\\
Sabiduría & \heb{חָכְמָה} & \emph{jojmá} & Sincronizador & Pr 3:19 \\
Verdadero & \heb{אֱמֶת} & \emph{emet} & Firmeza, realidad & Éx 34:6 \\
Misericordioso & \heb{רַחֲמִים} & \emph{rajamim} & Compasión entrañable & Sal 103:13 \\
Fiel & \heb{אֱמוּנָה} & \emph{emuná} & Constancia en el pacto & Dt 32:4 \\
Celo & \heb{קִנְאָה} & \emph{qin’á} & Protege su gloria y su pueblo & Éx 34:14 \\
Ira santa & \heb{אַף} & \emph{af} & Santidad ante el mal & Sal 2:12 \\
Bondad, paciencia, gracia & \heb{טוֹב}, \heb{אֶרֶךְ אַפַּיִם}, \heb{חֵן} & \emph{tov}, \emph{érej apáyim}, \emph{jen} & Favor y tolerancia & Éx 34:6; Sal 145:9 \\ \bottomrule
\end{tabular}
\par\smallskip{\footnotesize\color{sepia}\textit{Nota.} En Sal 115:3 el término \heb{אֲדֹנָי} nombra el concepto, pero no aparece en el verso.}
\end{table}

\section{Segunda clasificación: comunicables e incomunicables}

La tradición reformada ofrece una segunda forma de ordenar los atributos, complementaria de la anterior: según si Dios los comparte o no con el ser humano. Los \destacar{incomunicables} (no transferibles) son aquellos que solo Dios posee; los \destacar{comunicables} (transferibles) son aquellos que Dios posee en perfección y que refleja, en forma finita, en el ser humano creado a su imagen (Berkhof, 1938; Grudem, 2007).

\begin{verso}{Isaías 46:9}
Acordaos de las cosas pasadas desde los tiempos antiguos; porque yo soy Dios, y no hay otro Dios, y nada hay semejante a mí…
\end{verso}

\begin{filologia}{Is 46:9}
\heb{וְאֵין עוֹד} \emph{ve-en od} & \heb{אין} & “Y no hay otro”: exclusividad absoluta.\\
\heb{וְאֶפֶס כָּמוֹנִי} \emph{ve-éfes kamoní} & \heb{אפס} & “Y nada como yo”. \heb{אֶפֶס} significa “fin, cesación, nada”: no existe término de comparación. Es la base de los atributos incomunicables.\\
\end{filologia}

\begin{verso}{Génesis 1:27}
Y creó Dios al hombre a su imagen, a imagen de Dios lo creó; varón y hembra los creó.
\end{verso}

\begin{filologia}{Gn 1:27}
\heb{וַיִּבְרָא} \emph{vayivrá} & \heb{ברא} & “Y creó”: verbo reservado en la Biblia hebrea para la acción de Dios.\\
\heb{בְּצַלְמוֹ} \emph{be-tsalmó} & \heb{צלם} & “A su imagen”: representación, reflejo. Junto con \heb{דְּמוּת} (\emph{demut}, “semejanza”, Gn 1:26) es la base de los atributos comunicables.\\
\end{filologia}

Lo comunicable no es solo un privilegio; es también un mandato. Pedro cita Levítico para recordarlo:

\begin{verso}{1 Pedro 1:15–16}
sino, como aquel que os llamó es santo, sed también vosotros santos en toda vuestra manera de vivir; porque escrito está: Sed santos, porque yo soy santo.
\end{verso}

\begin{filologia}{1 P 1:15–16}
\grc{ἅγιος} \emph{hágios} & \grc{ἅγιος} & “Santo”: apartado para Dios. Traduce el hebreo \heb{קָדוֹשׁ}.\\
\grc{γενήθητε} \emph{genéthete} & \grc{γίνομαι} & “Sed, llegad a ser”: la santidad se va formando en la vida concreta.\\
\grc{γέγραπται} \emph{gégraptai} & \grc{γράφω} & “Escrito está”: cita de Lv 11:44, \heb{וִהְיִיתֶם קְדֹשִׁים} (“y seréis santos”).\\
\end{filologia}

La Tabla~\ref{tab:cruce} cruza las dos clasificaciones.

\begin{table}[htbp]
\centering\small
\caption{Cruce de las dos clasificaciones}\label{tab:cruce}
\renewcommand{\arraystretch}{1.3}
\begin{tabular}{@{}P{2.6cm} P{3.6cm} P{4.4cm} P{4.6cm}@{}}
\toprule
 & \textbf{Ser} & \textbf{Identidad} & \textbf{Relación} \\ \midrule
\textbf{Incomunicables} & Uno & Eterno, Inmutable, Omnipresente, Omnipotente*, Omnisciente* & — \\
\textbf{Comunicables} & Amor, Espíritu & Santo, Justo, Soberano (reflejado) & Sabiduría, Verdad, Misericordia, Fidelidad, Celo, Ira santa, Bondad, Paciencia, Gracia \\ \bottomrule
\end{tabular}
\par\smallskip{\footnotesize\color{sepia}\textit{Nota.} * Lo infinito no se comunica, pero sí su reflejo finito: poder y conocimiento. Siguiendo a Berkhof (1938), la espiritualidad se cuenta entre los comunicables: el ser humano recibe espíritu (Gn 2:7), no la espiritualidad infinita de Dios.}
\end{table}

\section{La sincronía de los atributos}

Si los atributos no se oponen, ¿qué los mantiene en armonía? En el esquema de esta obra, ese papel lo cumple la \destacar{sabiduría}: Dios aplica su justicia con misericordia, su poder con santidad y su amor con verdad, porque todo lo hace con sabiduría.

\begin{verso}{Proverbios 3:19}
\SENOR{} con sabiduría fundó la tierra; afirmó los cielos con inteligencia.
\end{verso}

\begin{filologia}{Pr 3:19}
\heb{בְּחָכְמָה} \emph{be-jojmá} & \heb{חכם} & “Con sabiduría”: destreza práctica, como la del artesano (Éx 31:3).\\
\heb{יָסַד} \emph{yasad} & \heb{יסד} & “Fundó”: echar los cimientos.\\
\heb{בִּתְבוּנָה} \emph{bi-tvuná} & \heb{בין} & “Con inteligencia”: capacidad de discernir y distinguir.\\
\end{filologia}

\begin{figure}[htbp]
\centering
%<fig:gloria_esquema>
\begin{tikzpicture}[font=\fontsize{10}{12}\selectfont,
  atSer/.style={draw=cSer, line width=0.5pt, fill=white, rounded corners=2pt, minimum width=2.3cm, minimum height=1.2cm, align=center, inner sep=2pt},
  hubSer/.style={atSer, line width=1.1pt, fill=cSer!7, minimum width=2.6cm},
  bandaSer/.style={draw=cSer!55, fill=cSer!3, line width=0.5pt, rounded corners=6pt},
  rotuloSer/.style={font=\fontsize{9}{11}\selectfont\scshape, text=cSer, anchor=north west},
  atIde/.style={draw=cIde, line width=0.5pt, fill=white, rounded corners=2pt, minimum width=2.3cm, minimum height=1.2cm, align=center, inner sep=2pt},
  hubIde/.style={atIde, line width=1.1pt, fill=cIde!7, minimum width=2.6cm},
  bandaIde/.style={draw=cIde!55, fill=cIde!3, line width=0.5pt, rounded corners=6pt},
  rotuloIde/.style={font=\fontsize{9}{11}\selectfont\scshape, text=cIde, anchor=north west},
  atRel/.style={draw=cRel, line width=0.5pt, fill=white, rounded corners=2pt, minimum width=2.3cm, minimum height=1.2cm, align=center, inner sep=2pt},
  hubRel/.style={atRel, line width=1.1pt, fill=cRel!7, minimum width=2.6cm},
  bandaRel/.style={draw=cRel!55, fill=cRel!3, line width=0.5pt, rounded corners=6pt},
  rotuloRel/.style={font=\fontsize{9}{11}\selectfont\scshape, text=cRel, anchor=north west},
  enl/.style={draw=gGris!70, line width=0.4pt},
  cruce/.style={draw=gGris, line width=0.45pt, densely dashed}]
\definecolor{gAzul}{RGB}{21,67,124}\definecolor{gAzulo}{RGB}{11,31,58}\definecolor{gGris}{RGB}{84,101,124}%
\definecolor{cSer}{RGB}{160,124,48}\definecolor{cIde}{RGB}{21,67,124}\definecolor{cRel}{RGB}{112,48,56}%
% Bandas
\draw[bandaSer] (-7.6,8.65) rectangle (7.6,11.35);
\draw[bandaIde] (-7.6,3.7) rectangle (7.6,8.4);
\draw[bandaRel] (-7.6,-1.05) rectangle (7.6,3.45);
\node[rotuloSer] at (-7.45,11.25) {I. Ser · ¿qué es?};
\node[rotuloIde] at (-7.45,8.3) {II. Identidad · ¿quién es?};
\node[rotuloRel] at (-7.45,3.35) {III. Relación · ¿cómo se revela?};
% I. Ser
\node[hubSer] (uno) at (0,9.85) {\textsc{Uno}\\[-1pt]{\fontsize{12}{13}\selectfont\heb{אֶחָד}}\\[-2pt]{\fontsize{8}{9.5}\selectfont\itshape ejad}};
\node[atSer] (amor) at (-3.6,9.85) {\textsc{Amor}\\[-1pt]{\fontsize{12}{13}\selectfont\heb{חֶסֶד}}\\[-2pt]{\fontsize{8}{9.5}\selectfont\itshape jésed}};
\node[atSer] (esp) at (3.6,9.85) {\textsc{Espíritu}\\[-1pt]{\fontsize{12}{13}\selectfont\heb{רוּחַ}}\\[-2pt]{\fontsize{8}{9.5}\selectfont\itshape rúaj}};
% II. Identidad
\node[hubIde] (santo) at (0,6.2) {\textsc{Santo}\\[-1pt]{\fontsize{12}{13}\selectfont\heb{קָדוֹשׁ}}\\[-2pt]{\fontsize{8}{9.5}\selectfont\itshape qadosh}};
\node[atIde] (eterno) at (-4.8,7.15) {\textsc{Eterno}\\[-1pt]{\fontsize{12}{13}\selectfont\heb{עוֹלָם}}\\[-2pt]{\fontsize{8}{9.5}\selectfont\itshape olam}};
\node[atIde] (sob) at (-4.8,5.8) {\textsc{Soberano}\\[-1pt]{\fontsize{12}{13}\selectfont\heb{אֲדֹנָי}}\\[-2pt]{\fontsize{8}{9.5}\selectfont\itshape Adonay}};
\node[atIde] (justo) at (-4.8,4.45) {\textsc{Justo}\\[-1pt]{\fontsize{12}{13}\selectfont\heb{צֶדֶק}}\\[-2pt]{\fontsize{8}{9.5}\selectfont\itshape tsedeq}};
\node[atIde] (omnis) at (4.8,7.15) {\textsc{Omnisciente}\\[-1pt]{\fontsize{12}{13}\selectfont\heb{יָדַע}}\\[-2pt]{\fontsize{8}{9.5}\selectfont\itshape yadá}};
\node[atIde, minimum width=2.6cm] (omnipr) at (4.8,5.8) {\textsc{Omnipresente}\\[-1pt]{\fontsize{12}{13}\selectfont\heb{מְלֹא כָל־הָאָרֶץ}}\\[-2pt]{\fontsize{8}{9.5}\selectfont\itshape meló kol-haárets}};
\node[atIde] (omnipo) at (4.8,4.45) {\textsc{Omnipotente}\\[-1pt]{\fontsize{12}{13}\selectfont\heb{אֵל שַׁדַּי}}\\[-2pt]{\fontsize{8}{9.5}\selectfont\itshape El Shadday}};
\node[atIde] (inm) at (0,4.45) {\textsc{Inmutable}\\[-1pt]{\fontsize{12}{13}\selectfont\heb{לֹא שָׁנִיתִי}}\\[-2pt]{\fontsize{8}{9.5}\selectfont\itshape lo shaniti}};
% III. Relación
\node[hubRel] (sab) at (0,1.15) {\textsc{Sabiduría}\\[-1pt]{\fontsize{12}{13}\selectfont\heb{חָכְמָה}}\\[-2pt]{\fontsize{8}{9.5}\selectfont\itshape jojmá}};
\node[atRel] (verd) at (-3.1,2.15) {\textsc{Verdadero}\\[-1pt]{\fontsize{12}{13}\selectfont\heb{אֱמֶת}}\\[-2pt]{\fontsize{8}{9.5}\selectfont\itshape emet}};
\node[atRel] (bond) at (-6.1,1.6) {\textsc{Bondad}\\[-1pt]{\fontsize{12}{13}\selectfont\heb{טוֹב}}\\[-2pt]{\fontsize{8}{9.5}\selectfont\itshape tov}};
\node[atRel] (ira) at (-6.1,-0.05) {\textsc{Ira santa}\\[-1pt]{\fontsize{12}{13}\selectfont\heb{אַף}}\\[-2pt]{\fontsize{8}{9.5}\selectfont\itshape af}};
\node[atRel] (celo) at (-3.1,-0.1) {\textsc{Celo}\\[-1pt]{\fontsize{12}{13}\selectfont\heb{קִנְאָה}}\\[-2pt]{\fontsize{8}{9.5}\selectfont\itshape qin'á}};
\node[atRel, minimum width=2.6cm] (mis) at (3.1,2.15) {\textsc{Misericordioso}\\[-1pt]{\fontsize{12}{13}\selectfont\heb{רַחֲמִים}}\\[-2pt]{\fontsize{8}{9.5}\selectfont\itshape rajamim}};
\node[atRel] (gra) at (6.1,1.6) {\textsc{Gracia}\\[-1pt]{\fontsize{12}{13}\selectfont\heb{חֵן}}\\[-2pt]{\fontsize{8}{9.5}\selectfont\itshape jen}};
\node[atRel] (pac) at (6.1,-0.05) {\textsc{Paciencia}\\[-1pt]{\fontsize{12}{13}\selectfont\heb{אֶרֶךְ אַפַּיִם}}\\[-2pt]{\fontsize{8}{9.5}\selectfont\itshape érej apáyim}};
\node[atRel] (fiel) at (3.1,-0.1) {\textsc{Fiel}\\[-1pt]{\fontsize{12}{13}\selectfont\heb{אֱמוּנָה}}\\[-2pt]{\fontsize{8}{9.5}\selectfont\itshape emuná}};
\node[font=\fontsize{8}{9.5}\selectfont\itshape, text=cRel] at (0,0.22) {sincroniza todos los atributos};
% Enlaces internos
\draw[enl] (amor) -- (uno) -- (esp);
\foreach \n in {eterno,sob,justo,omnis,omnipr,omnipo,inm} \draw[enl] (santo) -- (\n);
\foreach \n in {verd,bond,ira,celo,mis,gra,pac,fiel} \draw[enl] (sab) -- (\n);
% Enlaces entre las tres dimensiones
\draw[cruce] (amor.south) -- (eterno.north);
\draw[cruce] (esp.south) -- (omnis.north);
\draw[cruce] (justo.south) -- (verd.north west);
\draw[cruce] (santo.south west) to[out=-115,in=115] (sab.north west);
\draw[cruce] (santo.south east) to[out=-65,in=65] (sab.north east);
\draw[cruce] (ira.east) to[out=-10,in=190] (celo.west);
% Plenitud en Cristo
\draw[cSer, line width=0.9pt] (0,-1.55) -- (0,-2.35) (-0.27,-1.78) -- (0.27,-1.78);
\node[font=\fontsize{9}{11}\selectfont\scshape, text=gAzulo] at (0,-2.7) {Plenitud en Cristo · Col 2:9};
\end{tikzpicture}
%</fig:gloria_esquema>
\caption{Gloria de Dios: su esencia infinita}\label{fig:gloria}
\par{\footnotesize\color{sepia}\textit{Nota.} Esquema del autor: los atributos del Ser, de la Identidad y de la Relación, sincronizados por la Sabiduría.}
\end{figure}

\precision{asignar a la sabiduría el papel de “sincronizador” es una propuesta de esta obra, construida sobre textos como Pr 3:19 y Ro 11:33; no es una categoría de la teología clásica.}

\section{La plenitud en Cristo}

Cada capítulo de este libro terminará con un apartado titulado \destacar{“El atributo en Cristo”}. La razón es teológica: si el Nombre de Dios es la colección de sus atributos, Cristo es ese Nombre hecho carne. En Él los atributos no aparecen por partes, sino en plenitud.

\begin{verso}{Colosenses 1:19–20}
por cuanto agradó al Padre que en él habitase toda plenitud, y por medio de él reconciliar consigo todas las cosas, así las que están en la tierra como las que están en los cielos, haciendo la paz mediante la sangre de su cruz.
\end{verso}

\begin{filologia}{Col 1:19–20}
\grc{εὐδόκησεν} \emph{eudókesen} & \grc{δοκέω} & “Agradó”: fue la decisión complacida de Dios.\\
\grc{πᾶν τὸ πλήρωμα} \emph{pan to pléroma} & \grc{πληρόω} & “Toda la plenitud”: lo que llena por completo, sin que falte nada.\\
\grc{κατοικῆσαι} \emph{katoikésai} & \grc{οἰκέω} & “Habitar”: residir de forma permanente, no como visitante.\\
\grc{εἰρηνοποιήσας} \emph{eirenopoiésas} & \grc{εἰρήνη} & “Haciendo la paz”: en la cruz convergen la justicia y la misericordia (Sal 85:10).\\
\end{filologia}

\begin{verso}{Colosenses 2:9}
Porque en él habita corporalmente toda la plenitud de la Deidad,
\end{verso}

\begin{filologia}{Col 2:9}
\grc{θεότητος} \emph{theótetos} & \grc{θεός} & “De la Deidad”: la esencia divina misma, no solo cualidades divinas (\grc{θειότης}, Ro 1:20).\\
\grc{σωματικῶς} \emph{somatikós} & \grc{σῶμα} & “Corporalmente”: la plenitud de Dios en un cuerpo humano real.\\
\end{filologia}

\begin{verso}{Hebreos 1:3}
el cual, siendo el resplandor de su gloria, y la imagen misma de su sustancia, y quien sustenta todas las cosas con la palabra de su poder, habiendo efectuado la purificación de nuestros pecados por medio de sí mismo, se sentó a la diestra de la Majestad en las alturas,
\end{verso}

\begin{filologia}{He 1:3}
\grc{ἀπαύγασμα} \emph{apaúgasma} & \grc{αὐγή} & “Resplandor”: la luz que irradia de la fuente, inseparable de ella.\\
\grc{χαρακτήρ} \emph{charaktér} & \grc{χαράσσω} & “Imagen misma”: la impronta exacta que deja un sello sobre la cera.\\
\grc{ὑποστάσεως} \emph{hypostáseos} & \grc{ὑφίστημι} & “Sustancia”: la realidad subyacente del ser de Dios.\\
\end{filologia}

El evangelio de Juan muestra que en Cristo se cumple lo que Moisés pidió ver en el Sinaí:

\begin{verso}{Juan 1:14}
Y aquel Verbo fue hecho carne, y habitó entre nosotros (y vimos su gloria, gloria como del unigénito del Padre), lleno de gracia y de verdad.
\end{verso}

\begin{filologia}{Jn 1:14}
\grc{ἐσκήνωσεν} \emph{eskénosen} & \grc{σκηνή} & “Habitó”; literalmente “puso su tienda”. Evoca el tabernáculo (\heb{מִשְׁכָּן}, \emph{mishkán}) donde moraba la gloria.\\
\grc{ἐθεασάμεθα τὴν δόξαν} \emph{etheasámetha ten dóxan} & \grc{θεάομαι} & “Vimos su gloria”: lo que Moisés no pudo ver de frente (Éx 33:20).\\
\grc{χάριτος καὶ ἀληθείας} \emph{cháritos kai aletheías} & — & “Gracia y verdad”: eco ampliamente reconocido de \heb{חֶסֶד וֶאֱמֶת} de Éx 34:6.\\
\end{filologia}

\begin{verso}{Juan 14:9}
Jesús le dijo: ¿Tanto tiempo hace que estoy con vosotros, y no me has conocido, Felipe? El que me ha visto a mí, ha visto al Padre; ¿cómo, pues, dices tú: Muéstranos el Padre?
\end{verso}

\begin{filologia}{Jn 14:9}
\grc{ἔγνωκάς με} \emph{égnokás me} & \grc{γινώσκω} & “Me has conocido”: el conocimiento de Jer 9:24 se realiza en la persona de Jesús.\\
\grc{ὁ ἑωρακὼς ἐμέ} \emph{ho heorakós emé} & \grc{ὁράω} & “El que me ha visto”: perfecto griego, un ver con efecto permanente.\\
\end{filologia}

\begin{verso}{2 Corintios 4:6}
Porque Dios, que mandó que de las tinieblas resplandeciese la luz, es el que resplandeció en nuestros corazones, para iluminación del conocimiento de la gloria de Dios en la faz de Jesucristo.
\end{verso}

\begin{filologia}{2 Co 4:6}
\grc{φωτισμόν} \emph{fotismón} & \grc{φῶς} & “Iluminación”: Dios repite en el corazón lo que hizo en la creación (Gn 1:3).\\
\grc{γνώσεως τῆς δόξης} \emph{gnóseos tes dóxes} & \grc{γινώσκω} & “Conocimiento de la gloria”: une el conocer de Jer 9:24 con la gloria de Éx 33:18.\\
\grc{ἐν προσώπῳ} \emph{en prosópo} & \grc{πρόσωπον} & “En la faz”: lo que Moisés no vio en el rostro de Dios, se ve en el rostro de Cristo.\\
\end{filologia}

\section{Nuestra gloria}

El libro termina donde comenzó, porque el propio Nuevo Testamento cierra el círculo. Pablo cita Jeremías 9:24 y lo aplica a Cristo:

\begin{verso}{1 Corintios 1:30–31}
Mas por él estáis vosotros en Cristo Jesús, el cual nos ha sido hecho por Dios sabiduría, justificación, santificación y redención; para que, como está escrito: El que se gloría, gloríese en el Señor.
\end{verso}

\begin{filologia}{1 Co 1:30–31}
\grc{ὁ καυχώμενος} \emph{ho kaujómenos} & \grc{καυχάομαι} & “El que se gloría”. Es el verbo con que la Septuaginta traduce \heb{יִתְהַלֵּל} en Jer 9:22–23: Pablo cita deliberadamente a Jeremías.\\
\grc{σοφία} \emph{sofía} & \grc{σοφός} & “Sabiduría”: lo que el sabio de Jer 9:23 buscaba en sí mismo, ahora es una Persona.\\
\grc{ἐν κυρίῳ} \emph{en kyrío} & \grc{κύριος} & “En el Señor”: en la Septuaginta \grc{κύριος} traduce el Nombre \heb{יְהוָה}; Pablo lo aplica a Cristo.\\
\end{filologia}

La sabiduría, la fuerza y la riqueza en que el ser humano quería gloriarse, Dios las ha dado en Cristo. Por eso, conocer los atributos de Dios no es un ejercicio intelectual: es aprender a gloriarnos en el lugar correcto.

\section{Plan de la obra}

El libro se organiza en tres partes que siguen la primera clasificación, y una conclusión. Cada capítulo desarrollará un atributo con la misma estructura: (1) texto ancla en RVR1960; (2) apartado filológico de su raíz; (3) \destacar{el Nombre que lo revela}; (4) definición; (5) teología; (6) un encuentro en la historia de Israel; (7) \destacar{el atributo en Cristo}; y (8) aplicación a la vida.

\begin{table}[htbp]
\centering\small
\caption{Estructura del libro}\label{tab:plan}
\renewcommand{\arraystretch}{1.25}
\begin{tabular}{@{}P{3.4cm} P{12.2cm}@{}}
\toprule
\textbf{Parte} & \textbf{Capítulos} \\ \midrule
Introducción & Conocer a Dios en sus atributos \\
I. Ser & Preexistente · Uno · Amor · Espíritu \\
II. Identidad & Santo · Eterno · Soberano · Justo · Inmutable · Omnipotente · Omnipresente · Omnisciente \\
III. Relación & Sabiduría · Verdadero · Misericordioso · Fiel · Celo · Ira santa · Bondad, paciencia y gracia \\
Conclusión & La gloria de Dios: la comunión de sus atributos en Cristo \\ \bottomrule
\end{tabular}
\end{table}

\section{Conclusión}

Esta introducción ha establecido el punto de partida de toda la obra. Jeremías enseña que la única gloria legítima del ser humano es entender y conocer a Dios. Ese conocimiento no es una conquista humana, sino un don: Dios se da a conocer proclamando su Nombre, y su Nombre es la colección de sus atributos. Todos ellos son perfecciones eternas; unos expresan lo que Dios es en su Ser, otros quién es en su Identidad, y otros se fueron revelando en su Relación con el ser humano. Unos son exclusivamente suyos; otros los refleja en nosotros, creados a su imagen.

Pero la meta no es una lista. La \destacar{gloria de Dios es la comunión perfecta de todos sus atributos}, que actúan como uno solo porque Dios es uno. Y esa plenitud no quedó oculta: habita corporalmente en Cristo, en cuyo rostro brilla el conocimiento de la gloria de Dios. Por eso cada capítulo terminará en Él, y por eso el libro entero es una invitación a cumplir el mandato de Jeremías: \emph{el que se gloría, gloríese en el Señor}.

\section*{Referencias}\markboth{Referencias}{}
\refapa{Berkhof, L. (1938). \emph{Systematic theology}. Eerdmans.}
\refapa{Elliger, K., \& Rudolph, W. (Eds.). (1997). \emph{Biblia Hebraica Stuttgartensia} (5.ª ed.). Deutsche Bibelgesellschaft.}
\refapa{Grudem, W. (2007). \emph{Teología sistemática: Una introducción a la doctrina bíblica}. Editorial Vida.}
\refapa{Koehler, L., \& Baumgartner, W. (2001). \emph{The Hebrew and Aramaic lexicon of the Old Testament} (Ed. de estudio). Brill.}
\refapa{Nestle, E., \& Aland, K. (Eds.). (2012). \emph{Novum Testamentum Graece} (28.ª ed.). Deutsche Bibelgesellschaft.}
\refapa{Pink, A. W. (2020). \emph{Los atributos de Dios} (E. Rodríguez Cruz, Trad.). Chapel Library. (Obra original publicada en 1930)}
\refapa{Rahlfs, A., \& Hanhart, R. (Eds.). (2006). \emph{Septuaginta} (Ed. rev.). Deutsche Bibelgesellschaft.}
\refapa{\emph{Santa Biblia: Reina-Valera 1960}. (1960). Sociedades Bíblicas en América Latina.}

\vspace{6pt}
{\small\color{sepia}Las citas bíblicas se transcriben de la Reina-Valera 1960. Los términos hebreos siguen la BHS; los griegos, el NA28 y la Septuaginta de Rahlfs-Hanhart.}
\end{document}
